\section{Introduction}
\label{sec:intro}

A government that guarantees mortgages, lends for education and funds itself from payroll and
income taxes has taken a position on wages without ever describing it as one. If automation moves
income from labor to capital, that position loses value on both sides at once: the revenue that
services the debt falls, and the credit the government has guaranteed deteriorates for the same
reason. Whether it holds an offsetting claim on the capital that replaced the labor is an
empirical question, and it is the question this paper answers.

The answer is that it does not. Measured on the same basis and in the same period, the United
States holds the wage side of a two-sided bet on artificial intelligence and almost none of the
other side. That is a structural observation about who holds what, not a forecast about which
direction events will take, and it has a consequence in both directions: if AI succeeds and
displaces labor, wage taxes fall and wage-backed credit deteriorates; if AI disappoints, capital
gains and corporate receipts fall instead. The federal government loses in both cases, through
different tax bases, and in neither does it hold a claim large enough to offset the loss.

Three measurements carry the argument. About \result{DebtOnlyDirect} percent of all US debt,
public and private, is serviced directly from wages, and \result{DebtOnlyIndirect} percent once
indirect servicing is counted. The federal government is exposed, as holder, guarantor or debtor,
on about \result{SovereignUnion} percent of those directly wage-backed claims, against about
\result{AIFederalShareRounded} percent of the claims on AI capital. And income that moves from
wages to AI capital bears an effective marginal tax rate of about \result{TauKBarkai} percent,
where replacing the wage taxes lost would require \result{RequiredEasy} to \result{RequiredHard}
percent.

\paragraph{Contributions.} The paper makes four.

\begin{enumerate}[leftmargin=1.4em,itemsep=0.25em]
\item \textbf{A measurement framework.} Labor backing accounts decompose the whole stock of US
financial claims two ways at once, by the income that services each claim and by who ultimately
holds, guarantees or owes it. The framework yields a ratio, a federal exposure share, a holder map
for both sides of the exposure, a series running from 1952, and a breakdown by wage quintile. The
construction, the conventions, and the sensitivity of each result to the judgement calls behind it
are stated in full, including the one convention that does most of the work.

\item \textbf{A holder map that shows the asymmetry.} The wage side and the AI side are measured
on the same basis so that the two can be compared rather than asserted. The AI side is built from
on-balance-sheet filings alone, so it is a lower bound and the federal share of it an upper bound,
which means the asymmetry we report is the conservative version of it.

\item \textbf{An explanation of why the tax system cannot close the gap.} We rebuild the effective
marginal rate on income that shifts from wages to AI capital from components, taking the treatment
of expensing from \citet{AcemogluManeraRestrepo2020} rather than assuming it, and show that the
binding constraint is the ownership of corporate claims rather than offshore profit shifting.
About \result{ThetaUntaxable} percent of US corporate equity is held in forms the individual
income tax barely reaches, and about \result{GainsUntaxedAtDeath} percent of accrued gains are
never taxed because basis is stepped up at death.

\item \textbf{A measured incidence of displacement.} At the only displacement level inside the
observed data we measure where the first-round loss lands across the federal government,
households, the housing agencies and \result{NInstitutions} bank and credit union balance sheets,
and we test whether household relief protects bank capital. It does not, and the reason it does
not is measured rather than argued.
\end{enumerate}

\paragraph{What we do not claim.} We do not forecast the pace of automation, and we do not claim
that AI will displace any particular quantity of labor: the paper reports a response surface
rather than a point on it. Only one level, ten percent of the wage bill, sits inside the range of
the data used to calibrate reemployment; every figure above that is labelled scenario and reported
as a band. We do not claim the fiscal mechanism, which is established in the literature reviewed
in Section~\ref{sec:related} and which we take as our starting point rather than our finding. We
do not claim that the measurement framework is new in the strong sense: our search located no
prior construction of the same object, but that search was not systematic, so the claim is stated
as ``none located'' rather than ``none exists''. And we do not claim that the capital tax result
settles which tax base is the right one; we argue for a marginal, flow-based rate and state the
case against it plainly in Section~\ref{sec:fiscal}.

\paragraph{Research questions and objectives.} RQ1 asks what share of US financial claims is
serviced by labor income, and who holds those claims; RO1 constructs the labor backing accounts
and the holder map, in Section~\ref{sec:accounts}. RQ2 asks whether the capital tax system can
replace the revenue lost when income shifts from wages to AI capital; RO2 builds the effective
rate from components and compares it with the required rate, in Section~\ref{sec:fiscal}. RQ3 asks
where the loss from displacement lands, by holder; RO3 measures first-round incidence at the level
inside the data, in Section~\ref{sec:displacement}. RQ4 asks what the second round does to the
institutions that hold the credit, and whether household relief protects them; RO4 applies the
losses across individual balance sheets, in Section~\ref{sec:banks}. RQ5 asks what happens in the
opposite direction, if AI disappoints; RO5 prices that case on the same basis, in
Section~\ref{sec:bust}. RQ6 asks what follows for institutions; RO6 sets out the minimum
sufficient set of instruments at each regime, and what the evidence does not support, in
Section~\ref{sec:institutions}.

\paragraph{How to read the numbers.} Every quantitative statement in this paper carries a status
on its first use. \emph{Replicated} means rebuilt by an instance that did not write the code, from
raw data and a published brief alone, inside a tolerance stated before the rebuild.
\emph{Measured} means produced by the pipeline and clearing every applicable check, but not
independently rebuilt. \emph{Scenario} means arithmetic conditional on an assumption stated where
the number appears. Quantities that were withdrawn appear only in Section~\ref{sec:limitations},
which lists them with reasons. Every number in the text is generated from an artifact in the
replication package rather than typed, and the same artifacts generate the tables.
